\scanpage{144}
\mspage{2}
le diagramme
\diagram{\begin{tikzcd}[ampersand replacement=\&,column sep=small]
\deleted{\unsure{y|x}}x_1\arrow[rr,"u"]\arrow[dr,"\zeta"']\&\&x_0\\
\&\alpha(x')\arrow[ur,"\eta"']\&
\end{tikzcd}}
est \inserted{fbt} filtrante à g.

plus généralement, considérons
\diagram{\begin{tikzcd}[ampersand replacement=\&]
X'=C'\arrow[rr,"\alpha"]\arrow[dr,dash,"f'"']\&\&C=X\arrow[dl,dash,"f"]\\
\&S\&
\end{tikzcd}}
et posons nous la question quand la \unsure{flèche|formule} \unsure{fonctorielle}, par rapport à un préfaisceau $F$ \inserted{variable} sur $X$,
\[
 f_*(F)\longrightarrow f'_{\unsure{*|d}}(\alpha^*(F))
\]
est toujours un iso. (Le cas précédent est celui où $\deleted{\unsure{I}}S=X$, $f=\mathrm{id}_X$.) On trouve qu'il suffit que $\forall s_0\in\operatorname{Ob}\unsure{s_0|S_0}$, le foncteur
\[
 \alpha_{/s_0}:X'_{/s_0}\longrightarrow X_{/s_0}
\]
soit ``anodin pour les $\varprojlim$'', \deleted{et} pour ceci, il suffit que pour objet \inserted{$(x_0,u_0)$} de $X_{/s_0}$ (\unsure{où}
\[
 f(x_0)\xrightarrow{u_0}s_0
\]
) la catégorie $\,{}_{(x_0,u_0)\backslash}(X'_{/s_0})$ soit fbt. filtr. à g.
% archive sheet number: 144
